\documentclass{article}

\textwidth=6.5in
\textheight=8.5in
\voffset=-54pt
\oddsidemargin=0in
\evensidemargin=0in
\setlength{\parskip}{8pt}
\setlength{\parindent}{0pt}

\usepackage{workbook}

\usepackage[]{handout}

\begin{document}

\begin{center}
  \large\textsc{Problem Set 27 - Integration by Substitution}
  
      \pspace{12pt}
  \begin{problemonly}\large Name: \rule{5cm}{0.15mm}\\ \end{problemonly}
\end{center}

\begin{grading}
  \begin{enumerate}
    \item 10 points
    \item 5 points
    \item 10 points
    \item 8 points
    \item 12 points
    \item 5 points
\end{enumerate}
Total: 50 points.
  

  \textbf{One note:} Please be strict with students about using correct notation
  when dealing with the endpoints in definite integrals.  For instance,
  they can write $\displaystyle \int_2^3 2x \sin(x^2)\,dx = \int_4^9 \sin
  u\,du$, or they can write $\displaystyle \int_2^3 2x \sin (x^2)\,dx =
  \int_{{\color{red}x=\color{black}2}}^{{\color{red}x=\color{black}3}}
  \sin u\,du$, but writing $\displaystyle \int_2^3 2x \sin(x^2)\,dx =
  \int_2^3 \sin u\,du$ is incorrect, and you should take off a little bit
  (like half a point) every time this happens, so that students get used
  to the idea that we care about this.

  They can also choose to solve the indefinite integral first with substitution, and then plug in the bounds at the end. We want to discourage this though, because students tend to forget the bounds.
\end{grading}

For help with these problems, read pp. 787-796 in Gottlieb.
% If you'd like a little extra practice, Exercise 25.2 on
%   page 796 gives you some hints without giving all of the details.  You
%   can also try any problems from the workbook that you didn't get
%   to in class (solutions are posted on the course website).

\begin{enumerate}
\item 

\begin{grading}
10 points: 2/3/2/3.
\end{grading}




  Find or evaluate the following integrals using the given substitution.  Write a sentence or two explaining why the given substitution is a good choice. 
  
  Then, check your answers to {{{{\#1}}(a)}} and
  {{{{\#1}}(c)}} by differentiating them.

  \begin{enumerate}[topsep=0pt]
  \item
\begin{problem}[\vfill]
      $\displaystyle \int x \sin (x^2)\,dx$. Let $u = x^2$, 
      $du = \fitb{1in}{2x\,dx}$.
    \end{problem}

    \begin{solution}
      \sub{$u = x^2$} is a good substitution since $x^2$ is the inner
      function in the composition $\sin(x^2)$, and we see something
      related to $du = 2x\,dx$ in the integral, namely
      \subd{$\frac{1}{2}\,du = x\,dx$}.  Then,
      \begin{align*}
        \int \subd{x} \sin(\sub{x^2})\, \subd{dx} &= \int \sin(\sub{u}) \cdot
        \subd{\tfrac{1}{2}\ du} \\
        &= - \frac{1}{2} \cos u + C \\
        &= \boxed{- \frac{1}{2} \cos(x^2) + C}
      \end{align*}
      Let's check:
      \begin{align*}
        \frac{d}{dx} \left(-\frac{1}{2}\cdot \fcolorbox{red}{white}{$\displaystyle
            \cos \left(\fcolorbox{blue}{white}{$\displaystyle x^ {2}
                $}\right)$} + C \right)  
        = {} & -\frac{1}{2}\cdot \left(-\sin \left(x^ {2} \right)\cdot
          2 x\right) \\
        = {} & x \sin(x^2) \checkmark
      \end{align*}
    \end{solution}

  \item

    \note{If they forget to change the endpoints here and plug in
      $x$-endpoints for $u$, they'll get something undefined.  So, that
      will hopefully get them to think a little more about what's going
      on.} 

    \begin{problem}[\vfill]
      $\displaystyle \int_0^1 \frac{1}{(4x + 1)^3}\,dx$. Let $u =
      4x + 1$; $du=\fitb{1in}{4\,dx}$.
    \end{problem}

    \begin{solution}
      \sub{$u = 4x + 1$} is a good substitution since $4x + 1$ is the
      inner function in the composition $(4x + 1)^3$.  With this
      substitution, $du = 4\,dx$, so \subd{$\frac{1}{4}\, du = dx$}.
      Since $x$ goes from $0$ to $1$, $u = 4x + 1$ goes from $4(0) + 1 =
      1$ to $4(1) + 1 = 5$, so we can rewrite the original integral as
      \begin{align*}
        \int_0^1 \frac{1}{(\sub{4x + 1})^3}\,\subd{dx} &= \int_1^5
        \frac{1}{\sub{u}^3} \cdot \subd{\frac{1}{4}\, du} \\
        &= \frac{1}{4} \int_1^5 u^{-3}\, du \\
        &= \frac{1}{4} \left( \left. - \frac{1}{2} u^{-2} \right|_1^5
        \right) \\
        &= \left. - \frac{1}{8} \cdot \frac{1}{u^2} \right|_1^5 \\
        &= - \frac{1}{8} \left( \frac{1}{25} - 1 \right) \\
        &= - \frac{1}{8} \left( - \frac{24}{25} \right) \\
        &= \boxed{\frac{3}{25}}
      \end{align*}
      \emph{Note on notation:} If you prefer not to change the bounds
      of the integral to be in terms of $u$, your first line should
      instead read
      \[ \int_0^1 \frac{1}{({4x + 1})^3}\,{dx} = \int_{x=0}^{x=1}
      \frac{1}{u^3} \cdot \frac{1}{4}\, du.\] It is
      \color{red}\textbf{not} \color{black} correct to write
      \[ \int_0^1 \frac{1}{({4x + 1})^3}\,{dx} = \int_0^1
      \frac{1}{u^3} \cdot \frac{1}{4}\, du.\]
      (The latter notation says that $u$ is going from $0$ to $1$, which
      is not true.)
    \end{solution}


  \item
\begin{problem}[\vfill]
      $\displaystyle \int \frac{x}{x^2 + 9}\,dx$. Let $u = x^2 +
      9$; $du=\fitb{1in}{2x\,dx}$. 
    \end{problem}

    \begin{solution}
      \sub{$u = x^2 + 9$} looks like a reasonable substitution since it's
      in the denominator and we see something related to $du = 2x\,dx$ in
      the integral, namely \subd{$\frac{1}{2}\,du = x\,dx$}.  Then,
      \begin{align*}
        \int \frac{\subd{x}}{\sub{x^2 + 9}}\,\subd{dx} &= \int
        \frac{1}{u} \left( \frac{1}{2}\,du \right) \\
        &= \frac{1}{2} \int \frac{1}{u}\,du \\
        &= \frac{1}{2} \ln |u| + C \\
        &= \boxed{\frac{1}{2} \ln |x^2 + 9| + C}
      \end{align*}
      Let's check.  First, notice that $x^2 + 9$ is always positive,
      so $|x^2 + 9|$ is the same as $x^2 + 9$.  Therefore, another way
      of writing our answer is $\frac{1}{2} \ln (x^2 + 9) + C$:
      \[ \frac{d}{dx} \left(\frac{1}{2}
        \fcolorbox{red}{white}{$\displaystyle
          \ln\left(\fcolorbox{blue}{white}{$\displaystyle x^ { 2} +
              9$}\right)$}\right) = \frac{1}{2} \cdot \frac{1}{x^ { 2} +
        9}\cdot 2 x = \frac{x}{x^2 + 9} \checkmark\]
    \end{solution}

    \pspace{\newpage}

  \item
    \begin{problem}[\vfill]
      $\displaystyle \int_0^{\ln 4} e^{x} \sqrt{3 e^x+13}\,dx$.
      Let $u = 3 e^x+13$; $du=\fitb{1in}{3e^x\,dx}$.
    \end{problem}

    \begin{solution}
      \sub{$u = 3 e^x + 1$} looks like a good substitution since it's the
      inner function in the composition $\sqrt{3 e^x + 13}$, and we see
      something related to $du = 3 e^x\,dx$ in the integral, namely
      \subd{$\frac{du}{3} = e^x\,dx$}.  Since $x$ goes from $0$ to $\ln
      4$, $u$ goes from $3 e^0 + 13 = 16$ to $3 e^{\ln 4} + 13 = 3 \cdot 4
      + 13 = 25$:
      \begin{align*}
        \int_0^{\ln 4} \subd{e^x} \sqrt{\sub{3 e^x + 13}}\,\subd{dx} &=
        \int_{16}^{25} \sqrt{\sub{u}}\, \subd{\frac{du}{3}} \\
        &= \frac{1}{3} \int_{16}^{25} u^{1/2}\,du \\
        &= \left. \frac{1}{3} \cdot \frac{2}{3} u^{3/2} \right|_1^4 \\
        &= \frac{2}{9} \cdot 25^{3/2} - \frac{2}{9} \cdot 16^{3/2} \\
        &= \frac{2}{9} \cdot 5^3 - \frac{2}{9} \cdot 4^3 \\
        &= \boxed{\frac{122}{9}}
      \end{align*}
    \end{solution}

  \end{enumerate}

\item 

\begin{grading}
10 points. 5 for each part
\end{grading}

 In one of the following integrals, the substitution $u = e^{2x} + 1$
    works well; in the other, the substitution $u = e^x$ does (it's up to
    you to decide which is which).  Find both indefinite integrals, and check
    your answers by differentiating them.

      \begin{enumerate}
      \item
        \begin{problem}[\vfill]
          $\displaystyle \int \frac{2 e^{x}}{e^{2x} + 1}\,dx$
        \end{problem}

        \begin{solution} % Jason Bland
          Here, it works well to substitute \sub{$u = e^x$}, since we
          also see \subd{$du = e^x \, dx$} in the integral.  Note that
          $e^{2x}$ can be rewritten as $e^x \cdot e^x$, or $(e^x)^2$: 
          \begin{align*}
            \int \frac{2 \subd{e^x}}{(\sub{e^x})^2 + 1} \, \subd{dx} &=
            \int \frac{2}{\sub{u}^2 + 1} \, \subd{du} \\ 
            &= 2 \arctan u + C \\
            &= \boxed{2 \arctan(e^x) + C}
          \end{align*}
          Let's check:
          \[ \frac{d}{dx} \left(2 \fcolorbox{red}{white}{$\displaystyle
              \arctan\left(\fcolorbox{blue}{white}{$\displaystyle e^ {
                    x} $}\right)$}\right) 
          = \frac{2}{\left(e^ { x} \right)^2+1}\cdot e^ { x} \checkmark\]
        \end{solution}

      \item
        \begin{problem}[\vfill]
          $\displaystyle \int \frac{e^{2x}}{e^{2x} + 1}\,dx$
        \end{problem}

        \begin{solution} % Jason Bland
          Here, it works better to substitute \sub{$u = e^{2x} + 1$}
          because we see something related to $du = 2e^{2x} \, dx$ in
          the integral, namely \subd{$\frac{du}{2} = e^{2x} \, dx$}. So,
          \begin{align*}
            \int \frac{\subd{e^{2x}}}{\sub{e^{2x} + 1}} \, \subd{dx} &=
            \int \frac{1}{\sub{u}} \cdot \subd{\frac{du}{2}} \\
            &= \frac{1}{2} \int \frac{1}{u} \, du \\
            &= \frac{1}{2} \ln |u| + C \\
            &= \boxed{\frac{1}{2} \ln \left|e^{2x} + 1 \right| + C}
          \end{align*}
          Note that $e^{2x} + 1$ is always positive, so we can also
          rewrite this as simply $\frac{1}{2} \ln \left|e^{2x} + 1
          \right| + C$.  Let's check that this is right by
          differentiating it:
          \[ \frac{d}{dx} \left(\frac{1}{2}
            \fcolorbox{red}{white}{$\displaystyle
              \ln\left(\fcolorbox{blue}{white}{$\displaystyle e^ { 2 x}
                  + 1$}\right)$}\right) = \frac{1}{2}\cdot \frac{1}{e^ {
              2 x} + 1}\cdot 2 e^{2x} = \frac{e^{2x}}{e^{2x} + 1}
          \checkmark\]
        \end{solution}

      \end{enumerate}

      \pspace{\newpage}

\item 

\begin{grading}
5 points: 2/2/1.
\end{grading}

 Consider the integral $\displaystyle \int \cos x \sin x\,dx$.

    \begin{enumerate}[topsep=0pt]

    \item
      \begin{problem}[\vfill]
        Using the substitution $u=\sin x$, show that $\displaystyle \int
        \cos x \sin x\,dx=\frac{1}{2} \sin^2x +C$.
      \end{problem}

      \begin{solution} % Jason Bland
        With \sub{$u = \sin x$}, we have \subd{$du = \cos x \, dx$}, so
        \begin{align*}
          \int \subd{\cos x} \, \sub{\sin x} \, \subd{dx} &= \int
          \sub{u} \, \subd{du} \\ 
          &= \frac{1}{2} u^2 + C \\
          &= \frac{1}{2} \sin^2 x + C
        \end{align*}
      \end{solution}

    \item
      \begin{problem}[\vfill]
        Using the substitution $u=\cos x$, show that $\displaystyle \int
        \cos x \sin x\,dx=-\frac{1}{2} \cos^2x +C$.
      \end{problem}

      \begin{solution} % Jason Bland
        Using \sub{$u = \cos x$}, we have $du = -\sin x \, dx$, so
        \subd{$\sin x \, dx = -du$}. Therefore,
        \begin{align*}
          \int \sub{\cos x}\, \subd{\sin x \, dx} &= \int \sub{u} \,
          \subd{(-du)} \\
          &= - \int u\,du \\
          &= -\frac{1}{2} u^2 + C \\
          &= -\frac{1}{2} \cos^2 x + C
        \end{align*}
      \end{solution}

    \item
      \begin{problem}[\vfill]
        Show that, although these answers look different, they are
        really equivalent. \\ \emph{Hint: use the Pythagorean Identity.} 
      \end{problem}

      \begin{solution} % Jason Bland
        Remember that $\sin^2 x + \cos^2 x = 1$, so we can rewrite
        $\cos^2 x = 1 - \sin^2 x$.  Then,
        \[ -\frac{1}{2} \cos^2 x + C = -\frac{1}{2} (1 - \sin^2 x) + C =
        \frac{1}{2} \sin^2 x + \left( - \frac{1}{2} + C \right),\] so we
        see that $\displaystyle - \frac{1}{2} \cos^2 x + C$ is the same
        as $\dfrac{1}{2} \sin^2 x + \text{a constant}$.
      \end{solution}

    \end{enumerate}



\item 

\begin{grading}
8 points. 2 for each part.
\end{grading}

 Find the following indefinite integrals.  (It may
  take you a few tries to pick a useful $u$; that's normal when you're
  first learning substitution!)

  \probonly{}
    \begin{enumerate}
    \item
      \begin{problem}[\vfill]
        $\displaystyle \int \frac{1}{\sqrt{2x+1}}\,dx$
      \end{problem}

      \begin{solution}
        It looks like a good substitution is \sub{$u = 2x + 1$}, since $2x
        + 1$ is the inner function in the composition $\sqrt{2x + 1}$.
        With this substitution, $du = 2\,dx$, so \subd{$\frac{1}{2}\,du =
          dx$}.  Then,
        \begin{align*}
          \int \frac{1}{\sqrt{\sub{2x+1}}}\,\subd{dx} &= \int
          \frac{1}{\sqrt{u}} \cdot \subd{\frac{1}{2}\,du} \\
          &= \frac{1}{2} \int u^{-1/2}\,du \\
          &= u^{1/2} + C \\
          &= \boxed{\sqrt{2x + 1} + C}
        \end{align*}
      \end{solution}

      \pspace{\newpage}

    \item
      \begin{problem}[\vfill]
        $\displaystyle \int \frac{1}{2x+1}\,dx$
      \end{problem}

      \begin{solution}
        There's no composition here, but it could be helpful to substitute
        for the denominator, so let's try \sub{$u = 2x + 1$}.  With this
        substitution, $du = 2\,dx$, so \subd{$\frac{1}{2}\,du = dx$}.
        Then,
        \begin{align*}
          \int \frac{1}{\sub{2x+1}}\,\subd{dx} &= \int \frac{1}{\sub{u}}
          \cdot \subd{\frac{1}{2}\,du} \\
          &= \frac{1}{2} \ln |u| + C \\
          &= \boxed{\frac{1}{2} \ln |2x + 1| + C}
        \end{align*}
      \end{solution}

    \item
      \begin{problem}[\vfill]
        $\displaystyle \int x \sqrt{2x^2+1}\,dx$
      \end{problem}

      \begin{solution}
        It looks like a good substitution is \sub{$u = 2x^2 + 1$}, since
        $2x^2 + 1$ is the inner function in the composition $\sqrt{2x +
          1}$, and we see something related to $du = 4x\,dx$ in the
        integral, namely $\subd{\frac{du}{4} = x\,dx}$:
        \begin{align*}
          \int \subd{x} \sqrt{\sub{2x^2+1}}\,\subd{dx} &= \int
          \sqrt{\sub{u}} \subd{\frac{du}{4}} \\
          &= \frac{1}{4} \int u^{1/2}\,du \\
          &= \frac{1}{4} \cdot \frac{2}{3} u^{3/2} + C \\
          &= \boxed{\frac{1}{6} (2x^2 + 1)^{3/2} + C}
        \end{align*}
      \end{solution}

    \item
      \begin{problem}[\vfill]
        $\displaystyle \int \frac{x}{\sqrt{2x^2+1}}\,dx$
      \end{problem}

      \begin{solution}
        It looks like a good substitution is \sub{$u = 2x^2 + 1$}, since
        $2x^2 + 1$ is the inner function in the composition $\sqrt{2x +
          1}$, and we see something related to $du = 4x\,dx$ in the
        integral, namely $\subd{\frac{du}{4} = x\,dx}$:
        \begin{align*}
          \int \frac{\subd{x}}{\sqrt{\sub{2x^2 + 1}}}\,\subd{dx} &= \int
          \frac{1}{\sqrt{\sub{u}}} \cdot \subd{\frac{du}{4}} \\
          &= \frac{1}{4} \int u^{-1/2}\,du \\
          &= \frac{1}{4} \cdot 2 u^{1/2} + C \\
          &= \boxed{\frac{1}{2} \sqrt{2x^2 + 1} + C}
        \end{align*}
      \end{solution}

    \end{enumerate}    
  \probonly{}

\pspace{\newpage}

\item 

\begin{grading}
12 points. 3 for each part.
\end{grading}

 Find or evaluate the following
  integrals.  (Again, don't get discouraged if your first choice of $u$
  doesn't work out; with practice, you'll get better at picking $u$.)

  \probonly{}
    \begin{enumerate}
    \item
      \begin{problem}[\vfill]
        $\displaystyle \int_{1/4}^1 \frac{\cos \left( \sqrt{x}
          \right)}{\sqrt{x}}\,dx$
      \end{problem}

      \begin{solution} % Jason Bland
        A promising substitution is \sub{$u = \sqrt{x}$}, since it's the
        inner function in the composition $\cos \left( \sqrt{x} \right)$.
        Then, $du = \frac{1}{2} x^{-1/2}\,dx = \frac{1}{2\sqrt{x}}\,dx$,
        which is good because we see something related to this in the
        integral: \subd{$\frac{1}{\sqrt{x}} \, dx = 2 \, du$}.  As $x$
        goes from $\frac{1}{4}$ to $1$, $u = \sqrt{x}$ goes from
        $\sqrt{\frac{1}{4}} = \frac{1}{2}$ to $\sqrt{1} = 1$.  So,
        \begin{align*}
          \int_{1/4}^1 \frac{\cos (\sub{\sqrt{x}})}{\subd{\sqrt{x}}} \,
          \subd{dx}
          &= \int_{1/2}^1 \cos \sub{u} \cdot \subd{2\,du} \\
          &= 2 \int_{1/2}^1 \cos u \, du \\ 
          &= 2 \sin u \Big|_{1/2}^1 \\
          &= \boxed{2 \sin 1 - 2 \sin \frac{1}{2}}
        \end{align*}
      \end{solution}

    \item
      \begin{problem}[\vfill]
        $\displaystyle \int \sqrt{\cos x}\sin x \,dx$
      \end{problem}

      \begin{solution}
        It looks like a good substitution is \sub{$u = \cos x$}, since
        $\cos x$ is the inner function in the composition $\sqrt{\cos x}$;
        in addition, we see something related to $du = - \sin x\,dx$ in
        the integral, namely \subd{$-du = \sin x\,dx$}.  Then,
        \begin{align*}
          \int \sqrt{\cos x}\sin x \,dx &= \int \sqrt{\sub{u}}
          \left(\subd{-du} \right) \\
          &= - \int u^{1/2}\,du \\
          &= - \frac{2}{3} u^{3/2} + C \\
          &= \boxed{- \frac{2}{3} \left( \cos x \right)^{3/2} + C}
        \end{align*}
      \end{solution}

    \item
      \begin{problem}[\vfill\newpage]
        $\displaystyle \int \tan x\,dx$
      \end{problem}

      \begin{solution}
        Let's first rewrite $\tan x$ as $\displaystyle \frac{\sin x}{\cos
        x}$, so that we are trying to evaluate $\displaystyle \int
      \frac{\sin x}{\cos x}\,dx$.  It looks like \sub{$u = \cos x$} is a
      promising substitution because we also see something related to
      $du = - \sin x\,dx$, namely \subd{$-du = \sin x\,dx$}:
      \begin{align*}
        \int \frac{\subd{\sin x}}{\sub{\cos x}}\,\subd{dx} &= \int
        \frac{1}{\sub{u}} \left( \subd{-du} \right) \\
        &= - \int \frac{1}{u}\,du \\
        &= - \ln |u| + C \\
        &= \boxed{- \ln |\cos x| + C}
      \end{align*}
      \end{solution}

    \item
      \begin{problem}[\vfill]
        $\displaystyle \int_0^{\ln 5} \frac{3e^x}{\sqrt{e^x +4}}\, dx$
      \end{problem}

      \begin{solution}
        It looks like a good substitution is \sub{$u = e^x+4$}, since $e^x
        + 4$ is the inner function in the composition $\sqrt{e^x + 4}$; in
        addition, we see \subd{$du = e^x\,dx$} in the integral.  As $x$
        goes from $0$ to $\ln 5$, $u = e^x + 4$ goes from $e^0 + 4 = 5$ to
        $e^{\ln 5} + 4 = 9$:
        \begin{align*}
          \int_0^{\ln 5} \frac{3 \subd{e^x}}{\sqrt{\sub{e^x +4}}}\,
          \subd{dx} &= \int_5^9 \frac{3}{\sqrt{\sub{u}}}\,\subd{du} \\
          &= \int_5^9 3 u^{-1/2}\,du \\
          &= \left. 6 u^{1/2} \right|_5^9 \\
          &= \boxed{18 - 6 \sqrt{5}}
        \end{align*}
      \end{solution}

    \end{enumerate}      
  \probonly{}

\item 

\begin{grading}
5 points.
\end{grading}

 \textbf{Reflect Back.} 

  Four of the five following integrals turn into a constant multiple of $\displaystyle \int
  \sqrt{u}\,du$ with an appropriate substitution. Which ones? What is the substitution?
  \begin{center}
    \emph{Example:} Substituting $u = \sin x$ in $\displaystyle \int 2
    \sqrt{\sin x} \cdot \cos x\,dx$ gives $\displaystyle \int 2
    \sqrt{u}\,du$.
  \end{center}

  \probonly{}
    \begin{enumerate}
    \item
      \begin{problem}
        $\displaystyle \int \frac{\sqrt{\ln x}}{x}\,dx$      
      \end{problem}

      \begin{solution}
        Substituting $u=\ln x$ gives $\displaystyle\int \sqrt{\ln x} \cdot
        \frac{1}{x}\,dx = \int \sqrt{u}\,du$.
      \end{solution}

    \item 
      \begin{problem}
        $\displaystyle \int \sqrt{\tan x} \cdot \sec^2 x\,dx$
      \end{problem}

      \begin{solution}
        Substituting $u = \tan x$ gives $\displaystyle \int \sqrt{\tan x}
        \cdot \sec^2 x\,dx = \int \sqrt{u}\,du$.
      \end{solution}

    \item
      \begin{problem}
        $\displaystyle \int e^x \sqrt{e^{2x} + 1}\,dx$
      \end{problem}

      \begin{solution}
        No such substitution exists.  (You might try to substitute $u =
        e^{2x} + 1$, but then $du = 2 e^{2x}\,dx$.)
      \end{solution}

    \item
      \begin{problem}
        $\displaystyle \int e^{2x} \sqrt{e^{2x} + 1}\,dx$
      \end{problem}

      \begin{solution}
        Substituting $u = e^{2x} + 1$ gives $\displaystyle \int e^{2x}
        \sqrt{e^{2x} + 1}\,dx = \int \frac{1}{2} \sqrt{u}\,du$.
      \end{solution}

    \item
      \begin{problem}
        $\displaystyle \int \sqrt{3x + 1}\,dx$
      \end{problem}

      \begin{solution}
        Substituting $u = 3x + 1$ gives $\displaystyle \int \sqrt{3x +
          1}\,dx = \int \frac{1}{3} \sqrt{u}\,du$.
      \end{solution}

    \end{enumerate}    
    \probonly{}

\end{enumerate}

\spaceonly{\psreflection{
Optional reading for next class is \S 25.3 in Gottlieb.
}}

\end{document}